% 原文 pp.137–143（印刷 pp.1255–1261）
\section{Lecture 10：Wilson 圈、't Hooft 圈与 't Hooft 的禁闭图像}

\begin{keytakeaways}
本讲把 Wilson 圈（Wilson loop）与 't Hooft 圈（'t Hooft loop）作为一对互为对偶的非局域算符引进来：前者是电荷沿闭曲线 $C$ 的和乐，后者则被定义为「沿 $C$ 带奇性的联络」上的路径积分。
二者通过四维 $U(1)$ 理论的 Abel 对偶互换，对偶把复耦合送到 $-1/\tau$。
核心数学工具是：(10.1)–(10.3) 那一串对 2-形式 $G$ 与对偶联络 $B$ 的 Gaussian 积分；用 $\mathbb{Z}$-覆叠上的单值性实现的「非法规范变换」$T_{\gamma'}(C')\Psi(A)=\Psi(A^\phi)$；以及群 $\pi_1(G)^*\times\pi_1(G)$ 上的自然辛形式。
主要结论有：$W_\gamma(C)$ 与 $T_{\gamma'}(C')$ 满足 't Hooft 的相位对易关系 (10.4)；在 Higgs 相中 't Hooft 圈满足面积律，其证明归结为「$\phi$ 必须在该丛第一 Chern 类的 Poincaré 对偶曲面上取零」，于是作用量正比于张成曲面的极小面积（这同时解释了 Meissner 效应与 Taubes 对 Seiberg–Witten 不变量的分析）。
据此可把有质量相分为 Coulomb、Higgs、禁闭三类，每一类对应 $\pi_1(G)^*\times\pi_1(G)$ 的某个极大迷向子群；而 't Hooft 的禁闭图像并未解释禁闭，只是把它重新表述为：必须理解那种交换 $W$ 与 $T$ 的非线性对偶。
\end{keytakeaways}

% ===== 原文 p.137 =====
\subsection{'t Hooft 圈算符}
% 原文 §10.1：'t Hooft loop operator

让我们首先回顾四维中的 Abel 对偶（abelian duality），这是我们在第 8 讲讨论过的。\footnote{原文页脚注：\emph{Notes by Pavel Etingof and David Kazhdan}——本讲笔记由 Pavel Etingof 与 David Kazhdan 整理。}
考虑一个带 $\theta$ 角的自由 $U(1)$ 规范理论。
于是我们有两个无量纲耦合 $e,\theta$，它们组合成单个复耦合
\[
\tau=\frac{2\pi i}{e^2}+\frac{\theta}{\pi},
\]
而拉格朗日量为
\[
\mathcal{L}=\int_M\Bigl(\frac{i\bar\tau}{4\pi}F_+^2-\frac{i\tau}{4\pi}F_-^2\Bigr),
\]
其中 $F$ 是 $U(1)$ 联络 $A$ 的曲率，$F_+,F_-$ 是曲率的自对偶（selfdual）与反自对偶（antiselfdual）部分。
我们已经看到，若时空 $M$ 是 spin 流形，则该理论作为 $\tau$ 的函数是「模不变」（modular invariant）的。
一个模对称 $\tau\to\tau+1$ 是显然的，因为它对应于把 $\theta$ 角平移 $\pi$；当 $M$ 是 spin 流形时这不改变任何东西，因为 spin 流形的 $c_1^2$ 是偶数。
（在非 spin 流形上，该对称性只是 $\tau\to\tau+2$。）
模群第二个生成元下的对称性 $\tau\to-1/\tau$ 更有意思，它对应于十九世纪经典地发现的电磁对偶（electromagnetic duality）。
更确切地说，这意味着：耦合常数为 $\tau$ 的联络 $A$ 的理论，在经典与量子两个层面上，都等同于耦合为 $-1/\tau$、联络为 $B$（且 $dA=\mathrm{const}\ast dB$）的同一理论。
现在，像第 8 讲那样，我们想知道算符在这种对偶下会变成什么。
特别地，我们想知道 Wilson 圈算符会变成什么。

% [Tutor Note: 注意 $\bar\tau$ 配给 $F_+^2$、$\tau$ 配给 $F_-^2$，这与不少教材「$\tau$ 配 $F_+^2$」的写法相反，而这一配对是被规范动能项的正定性逼出来的。取 $F_\pm^2\equiv\int F_\pm\wedge\ast F_\pm$、$F_\pm=\tfrac12(F\pm\ast F)$，则 $F_+^2=\tfrac12(I_2+I_1)$、$F_-^2=\tfrac12(I_2-I_1)$，其中 $I_2=\int F\wedge\ast F$、$I_1=\int F\wedge F$；代入后得 $\mathcal{L}=\frac{1}{4e^2}\int F_{\mu\nu}F^{\mu\nu}+i\theta\,c_1^2$，$c_1^2=\frac{1}{4\pi^2}\int F\wedge F$。若把 $\tau$ 与 $\bar\tau$ 对调，$I_2$ 的系数会变成 $-\frac{1}{2e^2}$，规范动能项直接变负，而不只是 $\theta$ 项反号。本讲 (10.1)–(10.3) 的系数都与这一配对绑定：若照常见写法把 $\frac{i\tau}{4\pi}$ 配给 $F_+^2$，(10.2) 中的 $-i\bar\tau\int G_+^2$ 就会变成 $+i\tau\int G_+^2$，末了 (10.3) 的对偶耦合会从 $-1/\tau$ 变成 $-1/\bar\tau$，'t Hooft 圈的磁通与电荷的对应随之整体反号。]

% [Tutor Note: 「spin 流形的 $c_1^2$ 是偶数」是四维 spin 拓扑的标准事实，作者未给证明。依据是 Wu 公式：对任意 $x\in H^2(M;\mathbb{Z})$ 有 $x\cdot x\equiv x\cdot w_2(M)\pmod 2$，而 spin 条件恰是 $w_2(M)=0$，故相交形式是偶的，$c_1^2$ 自然为偶数（等价说法：spin 四维流形的相交形式是偶幺模格，这也是 Rokhlin 定理的出发点）。这也说明非 spin 流形上为何只剩 $\tau\to\tau+2$：此时 $c_1^2$ 可为奇数，$\theta\to\theta+\pi$ 会带来 $e^{i\pi c_1^2}=-1$ 的相位。参见 Milnor–Stasheff, \textit{Characteristic Classes}。]

回忆 Wilson 圈算符具有形式 $W_\gamma(C)=e^{i\gamma\,\mathrm{Hol}_C(A)}$，其中 $\mathrm{Hol}_C(A)$ 表示联络 $A$ 沿时空中一条闭定向曲线 $C$ 的积分。
当 $\gamma$ 为整数时（或者当曲线 $C$ 在 $M$ 中同伦平凡时，对任意实数 $\gamma$），该算符是规范不变的、且良定义。
（更一般地，可以有若干个分支 $C_i$ 与实数 $\gamma_i$，条件是 $\sum_i\gamma_iC_i$ 应当是 $H^1(M)$ 中的一个整类（integral class）。）
该算符的矩阵元照例通过把上述指数插入路径积分来计算。

% ===== 原文 p.138 =====
与我们在二、三维的类似问题中所得到的一样，我们应当得到：Wilson 圈的对偶描述是这样一种配方——它说，与其在路径积分中插入一个「生活在 $C$ 上」的对象，我们应当对沿 $C$ 具有奇性的联络作积分。

确切的答案是下面这个。
对任意曲线 $C$，期望值 $\langle W_\gamma(C)\mathcal{O}_1\cdots\mathcal{O}_n\rangle$ 等于 $\int e^{-L(B)}\mathcal{O}'_1\cdots\mathcal{O}'_n\,DB$，其中积分取遍 $M\setminus C$ 上某个丛的联络 $B$，使得 $B$ 的曲率在「垂直于 $C$、过任一点的三维空间中一个小 2-球面 $S$」上的积分等于 $2\pi\gamma$。
（这里 $\mathcal{O}_i$ 是任意一组位于与 $C$ 不相交的点处的局域算符；它们在对偶下变成局域算符 $\mathcal{O}'_i$。）
显然，$B$ 并不完全是 $U(1)$ 丛上的联络，但在数学上考虑这类联络并无困难。

让我们证明上一段的论断。
我们假定 $C$ 是一个边界。
设 $D$ 是以 $C$ 为边界的 2-链。
回忆第 8 讲中关于 Abel 对偶的计算：我们的场是 $A$——原来的联络，$G$——2-形式，以及 $B$——对偶联络。
我们有
\begin{equation}\tag{10.1}\label{eq:10-1}
\begin{aligned}
&\int e^{-L(A)}W_\gamma(C)\,DA\\
&\qquad=\int DA\,DG\,DB\,\exp\Bigl(\frac{-i\bar\tau}{4\pi}\int\mathcal{F}_+^2+\frac{i\tau}{4\pi}\int\mathcal{F}_-^2+\frac{i}{2\pi}\int G\wedge F_B+i\gamma\int_D(F_A-G)\Bigr),
\end{aligned}
\end{equation}
其中 $\mathcal{F}=F_A-G$，而最后一个因子对应于 Wilson 圈的插入（回忆 Abel 对偶的讨论：Wilson 圈经典地就是 $e^{i\gamma\int_D F_A}$，而在扩展理论中 $F_A$ 应代之以 $F_A-G$）。
把 $A$ 规范到 0，我们得到
\begin{equation}\tag{10.2}\label{eq:10-2}
\begin{aligned}
&\int e^{-L(A)}W_\gamma(C)\,DA\\
&\qquad=\int DG\,DB\,\exp\Bigl(\frac{-i\bar\tau}{4\pi}\int G_+^2+\frac{i\tau}{4\pi}\int G_-^2+\frac{i}{2\pi}\int G\wedge F_B-i\gamma\int_D G\Bigr)\\
&\qquad=\int DG\,DB\,\exp\Bigl(\frac{-i\bar\tau}{4\pi}\int G_+^2+\frac{i\tau}{4\pi}\int G_-^2+i\int G\wedge\bigl(\tfrac{1}{2\pi}F_B-\gamma[D]\bigr)\Bigr),
\end{aligned}
\end{equation}
其中 $[D]$ 是 $D$ 的 delta 函数。
现在定义 $\tilde F_B:=F_B-2\pi\gamma[D]$。
于是把 $G$ 积掉之后，(10.2) 可以写成
\begin{equation}\tag{10.3}\label{eq:10-3}
\int e^{-L(A)}W_\gamma(C)\,DA=\int DB\,\exp\Bigl(\frac{i}{4\pi\bar\tau}\int(\tilde F_B)_+^2-\frac{i}{4\pi\tau}\int(\tilde F_B)_-^2\Bigr),
\end{equation}

% [Tutor Note: (10.1)$\to$(10.3) 的两步作者都只用一句话带过，实际各有内容。①「把 $A$ 规范到 0」：这里其实没有可用来定 $A=0$ 的规范自由度，作者的意思是 $A$ 只以 $dA$ 的形式出现在 $\mathcal{F}=F_A-G$ 与 Wilson 因子 $i\gamma\int_D(F_A-G)$ 中，故可把 $G$ 平移 $G\to G+F_A$ 把它消掉（$DG$ 的 Lebesgue 测度在平移下不变），此时 $\mathcal{F}\to-G$、Wilson 因子变成 $-i\gamma\int_DG$，正好给出 (10.2)；等价地也可先对 $A$ 积分得到 $G=F_A$ 再代回。②「把 $G$ 积掉」：$G$ 的二次型 $-\frac{i\bar\tau}{4\pi}G_+^2+\frac{i\tau}{4\pi}G_-^2$ 在 $\mathrm{Im}\,\tau\neq0$ 时非退化，Gaussian 积分取系数之倒数，这正是 (10.3) 中 $\tau\to-1/\tau$ 的来源；线性项 $\frac{i}{2\pi}\int G\wedge F_B-i\gamma\int_D G=i\int G\wedge\bigl(\frac{1}{2\pi}F_B-\gamma[D]\bigr)$ 只起配方平移的作用，把 $F_B$ 换成 $\tilde F_B=F_B-2\pi\gamma[D]$。注意作为流有 $d\tilde F_B=-2\pi\gamma[C]$（符号依定向约定），但在 $M\setminus C$ 上 $\tilde F_B$ 是闭的，$\int_{S^2}\tilde F_B=2\pi\gamma$ 才是「$B$ 沿 $C$ 带奇性」的确切含义。]

因此，插入 $W_\gamma(C)$ 的效果是：在最后的答案中 $F_B$ 被 $\tilde F_B$ 取代。
这等价于说，$B$ 现在是沿 $C$ 带奇性的线丛上的联络，正如上面所断言的那样。
我们原先用「曲率在小 2-球面 $S$ 上的积分为 $2\pi\gamma$」给出的描述，指的就是曲率 $\tilde F_B$。

% ===== 原文 p.139 =====
设 $C$ 不是边界。
那么为了使 $W_\gamma(C)$ 有定义，我们必须取 $\gamma$ 为整数。
对这样的 $C$ 与 $\gamma$，$\langle W_\gamma(C)\rangle=0$（即使插入任意多个局域算符也是如此）。
的确，我们有一个对称性 $A\to A+d\phi$，而 $\phi$ 不必整体地定义为到圆周的映射；事实上，$d\phi$ 可以是任意闭 1-形式（更一般地，若存在挠率，$d\phi$ 可以是线丛上的任意平坦联络）。
因此，若 $C$ 不是边界，我们就可以这样选取 $d\phi$，使得算符 $W_\gamma(C)$ 乘以某个 $e^{i\alpha}$，其中 $\alpha\notin2\pi\mathbb{Z}$。
于是它的期望值（即使包含规范不变的局域算符，而这些算符在 $A\to A+d\phi$ 下不变）为零。
更一般地，关联函数 $\langle W_{\gamma_1}(C_1)\cdots W_{\gamma_n}(C_n)\rangle$（可带任意局域算符）为零，只要 $\sum\gamma_iC_i$ 在 $H_1(M)$ 中非零，其中 $M$ 是时空。
正如前面所注意到的，为使所讨论的这些算符之积良定义，我们只要求 $\sum\gamma_iC_i$ 是一个整类。
若该和是非零的整类，则算符之积良定义，而其期望值为零。

由我们迄今的构造，对任意曲线 $C$ 我们有两个算符：
\begin{enumerate}
\item $W_\gamma(C)=e^{i\gamma\int_C A}$；
\item $T_\gamma(C)=e^{i\gamma\int_C B}$。
\end{enumerate}
$B$ 是对偶理论的规范场。
第二个算符与 Wilson 圈对偶，称为 't Hooft 圈算符（'t Hooft loop operator）。
我们迄今所描述的，是如何用 $B$ 来表达 $W_\gamma$，或者如何用 $A$ 来表达 $T_\gamma$。

% [Tutor Note: 段首那句「$\sum\gamma_iC_i$ 在 $H_1(M)$ 中非零」与上一页的「$H^1(M)$ 中的整类」写法不同，但并不矛盾：$\gamma_iC_i$ 是 $M$ 中的实 1-链，「是整类」意即它的类落在 $H_1(M;\mathbb{Z})$ 的像中，而曲线只通过 $\int_C A$ 与联络配对，故用 $H_1$ 的语言或用其 Poincaré 对偶类的语言表述是同一条件。具体验证可取 $M=S^1\times S^2$、$C=S^1\times\{\mathrm{pt}\}$、$\gamma$ 非整数：$C$ 不同调于零，于是存在沿 $C$ 有非零周期的闭 1-形式 $d\phi$，$W_\gamma$ 乘上 $e^{i\gamma\oint_Cd\phi}\neq1$，而规范不变的局域算符不受影响，故 $\langle W_\gamma(C)\rangle=0$。]

\subsection{'t Hooft 圈算符的 Hilbert 空间解释}
% 原文 §10.2：Hilbert space interpretation of the 't Hooft loop operator

现在让我们从 Hamilton 的观点来看这幅图像。
此时时空具有形式 $M=M^3\times\mathbb{R}$，配以 Minkowski 度量。
设 $C,C'$ 是 $M^3$ 中两条不相交的闭简单曲线。
它们在 Hilbert 空间 $\mathcal{H}$ 上定义了算符 $W_\gamma(C)$ 与 $T_{\gamma'}(C')$。
我们把 $C$ 与 $C'$ 的环绕数（linking number）记为 $l(C,C')$。
下列对易关系归功于 't Hooft：
\begin{equation}\tag{10.4}\label{eq:10-4}
W_\gamma(C)T_{\gamma'}(C')=\exp\bigl(2\pi i\gamma\gamma'\,l(C,C')\bigr)T_{\gamma'}(C')W_\gamma(C).
\end{equation}

% [Tutor Note: (10.4) 的相位可以一行验出，也看得出符号从何而来：在 $\Psi(A)$ 表示中 $W_\gamma$ 是乘以 $e^{i\gamma\int_C A}$ 的算符，而 $T_{\gamma'}(C')$ 把 $A$ 变成 $A+d\phi$，其中 $\phi$ 绕 $C'$ 的单值性是 $e^{i\gamma'}$，即 $\oint_{S}d\phi=2\pi\gamma'$（$S$ 为与 $C'$ 连接的小圆）。于是 $T^{-1}WT=e^{i\gamma\oint_C d\phi}W$，而 $\oint_C d\phi=2\pi\gamma'\,l(C,C')$ 正是环绕数的定义，相位即 $e^{2\pi i\gamma\gamma' l}$。若同时反转 $C$ 与 $C'$ 的定向，$l\to-l$，对易子变成它的逆——这并非理论的不自洽，而只是「磁通方向」与「电荷符号」的相对约定。]

让我们证明这个公式。
我们按原来的联络 $A$ 来工作。
那么 Hilbert 空间由波函数 $\Psi(A)$ 构成。
在这一实现中，Wilson 圈算符 $W_\gamma(A)$ 就是乘以 $e^{i\gamma\int A}$ 的算符。

然而，'t Hooft 圈算符要难定义一些。
为此，考虑由与 $C'$ 的环绕数给出的同态 $\pi_1(M^3\setminus C',x_0)\to\mathbb{Z}$。
设 $\pi_1^0$ 为该同态的核，$X$ 为 $M_3\setminus C'$ 的、对应于 $\pi_1^0$ 的 $\mathbb{Z}$-覆叠。
设 $\phi$ 是函数 $X\to U(1)$，使得对应于 $\mathbb{Z}$ 的生成元的单值性（monodromy）为 $e^{i\gamma}$。
任意两个这样的函数相差一个规范变换，但 $\phi$ 本身并不是一个真正的（honest）规范变换。
于是不难验证，'t Hooft 圈算符 $T_{\gamma'}(C')$ 正是由 $\phi$ 所作的「非法」规范变换：
\begin{equation}\tag{10.5}\label{eq:10-5}
\bigl(T_{\gamma'}(C')\Psi\bigr)(A)=\Psi(A^\phi).
\end{equation}

% [Tutor Note: 原文此处似乎有笔误：既然 (10.5) 定义的是 $T_{\gamma'}(C')$，单值性应当是 $e^{i\gamma'}$ 而非 $e^{i\gamma}$；若照字面取 $e^{i\gamma}$，得到的将是 $T_\gamma(C')$，它与 (10.4) 的相位 $2\pi i\gamma\gamma'l$、以及下文「Wilson 圈由 $\pi_1(G)^*$ 分类、't Hooft 圈由 $\pi_1(G)$ 分类」都不相容。同页稍上方的 $M_3\setminus C'$ 也应是 $M^3\setminus C'$（全篇其余各处均写 $M^3$）。]

注意这是良定义的，因为任意两个这样的 $\phi$ 相差一个真正的规范变换。
为验证 (10.5) 成立，人们用 $M$ 上的路径积分来表示用 Hamilton 语言在 (10.5) 中表述的操作，然后把该路径积分与此前用「在连接 $C$ 的 2-球面 $S$ 上具有非平凡限制的联络」给出的 $T_\gamma(C)$ 的定义相比较。

% ===== 原文 p.140 =====
现在公式 (10.4) 就清楚了，因为 $TWT^{-1}=e^{i\gamma\int_C d\phi}W$，这是由于 $T$ 在 $W$ 的定义中变换联络的方式。

公式 (10.4) 可以推广到非 Abel 的 $G$（且 $\pi_1(G)\neq0$）。
为此，需要用沿 $C$ 带奇性的联络给出 $T_\gamma(C)$ 的定义；这个定义在 §10.3 中给出。
除此之外 (10.4) 的证明类似。

\subsection{禁闭的图像}
% 原文 §10.3：The picture of confinement

现在我们讨论 't Hooft 发展出来的一幅禁闭图像。
一般地，我们并不假定规范群是 Abel 的。
回忆禁闭的定义。
我们有规范群 $G$，以及它的万有覆盖（universal cover）$\tilde G$。
我们假定 $G$ 是 $\tilde G$ 的商，且 $G$ 忠实地作用在拉格朗日量中的所有场上。
设 $R$ 是 $\tilde G$ 的一个表示。
如第 7 讲所讨论，若存在质量间隙（mass gap），则 Wilson 圈算符在表示 $R$ 中的期望值 $\langle W_R(C)\rangle$ 随 $C$ 变大有两种常见的衰减模式：
模式 1：
\begin{equation}\tag{10.8}\label{eq:10-8}
\langle W_R(C)\rangle\sim e^{-\lambda\,\mathrm{Length}(C)}
\end{equation}
模式 2：
\begin{equation}\tag{10.9}\label{eq:10-9}
\langle W_R(C)\rangle\sim e^{-\lambda\,\mathrm{Area}(C)}.
\end{equation}
（这里 $\lambda$ 是某个固定的非零数，而 $C$ 的面积指张成曲面（spanning surface）的最小面积。）
第一种区域称为 Higgs 区（长度律），第二种称为禁闭区（面积律）。

% [Tutor Note: 公式编号在此从 (10.5) 直接跳到 (10.8)：本讲后面还有 (10.9)–(10.11)，而 (10.6)、(10.7) 全篇并未出现。合理的解释是上页「对任意曲线 $C$ 我们有两个算符」的 1)、2) 两条原本编号为 (10.6)、(10.7)，后来改成了列表却没有重排后续编号。译文按原文保留 (10.8)、(10.9) 的号码，读者不必以为自己漏读了两式。]

如第 7 讲所讨论，当 $R$ 是 $G$ 自身的表示时，总会得到模式 1。
有意思的情形是：$R$ 是 $\tilde G$ 的表示，却不是 $G$ 的表示。
类似地，若商 $\tilde G/G$ 以同样的方式作用在 $R$ 上，则两个 $R$ 可以认为是等价的。
于是，模去这个等价关系后，有意思的 $W_R(C)$ 由 $\pi_1(G)^*$ 分类。

现在让我们考虑 't Hooft 圈算符 $T_\gamma(C)$。
它对任意 $G$ 都可按照与 Abel 情形类比的方式如下定义。
我们固定一个元素 $\gamma\in\pi_1(G)$。
回忆：$S^2$ 上的 $G$-丛由一个取值于 $H^2(S^2,\pi_1(G))$ 的特征类分类，而后者典范同构于 $\pi_1(G)$。
因此 $\gamma$ 的选择典范地决定了 $S^2$ 上 $G$-丛的同构类。
现在我们可以定义 't Hooft 算符：带 $T_\gamma(C)$ 插入的路径积分，是对 $M\setminus C$ 上的联络作积分，这些联络具有如下性质：当限制到连接 $C$ 的小球面 $S$ 上时，丛的特征类为 $\gamma$。

't Hooft 的想法是：考虑 $T_\gamma(C)$ 而不是 $W_R(C)$，并找出其期望值服从面积律的条件。
如他所证明的，这在某些 Higgs 理论中确实发生。
于是 't Hooft 提出（沿循 Nambu、Mandelstam 等人的更早想法）：禁闭会通过一种把 't Hooft 圈算符映为 Wilson 圈算符的对偶，与 Higgs 机制联系起来。
这并没有解释禁闭，但它把问题重新表述为：要把禁闭这一神秘现象归结为远为容易理解的 Higgs 现象，人们必须理解那种交换 't Hooft 圈与 Wilson 圈算符的非线性对偶。

% ===== 原文 p.141 =====
为了说明 Higgs 理论中 't Hooft 圈的面积律，我们考虑一个熟悉的例子：带一个荷电复标量 $\phi$（电荷为 1）的 $U(1)$ 规范理论。
$n$ 维中的拉格朗日量为
\begin{equation}\tag{10.10}\label{eq:10-10}
\int d^nx\Bigl(\frac{|F|^2}{4e^2}+|D_A\phi|^2+V(\phi\bar\phi)\Bigr),
\end{equation}
其中 $V$ 是（四次）势。
我们将研究 't Hooft 圈 $T_\gamma(C)$，其中 $\gamma\in\mathbb{Z}$。
于是 $A$ 是一个联络，而 $\phi$ 是 $M\setminus C$ 上一个厄米线丛的截面，使得当限制到连接 $C$ 的小球面时，它具有第一 Chern 类 $\gamma$。
我们将对两类 $V$ 计算 $\langle T_\gamma(C)\rangle$：
\begin{enumerate}
\item $V=\lambda(\phi\bar\phi+a^2)^2$；
\item $V=\lambda(\phi\bar\phi-a^2)^2$。
\end{enumerate}

这个理论在第 2 讲中考虑过。回忆那次考虑的结果。

情形 1。
在红外，该理论的行为如同一个自由的有质量场与一个自由的无质量规范理论之积。
特别地，没有质量间隙。
（因此上面模式 1 与模式 2 的分类并不适用。）
于是，对于大的 $C$，我们可以用自由理论来计算 $\langle T_\gamma(C)\rangle$。
但在自由理论中，该期望值与对偶理论中的 $\langle W_\gamma(C)\rangle$ 相同。
容易看出，由于电荷相互作用的 Coulomb 律，$W_\gamma$ 与 $T_\gamma$ 的期望值都按长度律行为。
以这种方式行为的理论据说处于「Coulomb 相」。

情形 2。
在红外，该理论发生规范对称性破缺与 Higgs 机制。
如第 8 讲所分析的，谱由一个有质量的矢量与一个实的标量组成。
所以有质量间隙。
人们不相信这个理论表现出禁闭，也就是说，人们相信它对 Wilson 圈表现出长度律。
这当然正是微扰论所算出的。

在情形 2 中，我们将证明 't Hooft 圈算符存在面积律，其原因在于 Higgs 机制。
这出于拓扑上的理由，下面加以解释。
其中一些想法应当从第 8.4 讲对 Polyakov 模型的讨论中熟悉。

作为热身，考虑闭时空 $M$（维数 $n\ge2$）以及具有非平凡 $c_1$ 的线丛 $\mathcal{L}$。
让我们考虑我们的理论在这个丛的截面上的路径积分。
我们假定 $M$ 的度量非常「大」，即所有特征距离与 $1/m$ 相比都很大，其中 $m$ 是理论的质量间隙。
设 $R$ 为 $M$ 的特征半径。
结果是：该积分中所有场组态的作用量必定非常大，当 $R$ 大时其量级为 $(mR)^2$；它被「一个与 $M$ 无关的常数乘以代表与 $c_1(\mathcal{L})$ Poincaré 对偶的圈的最小二维曲面的面积」所下界。

% [Tutor Note: 「量级为 $(mR)^2$」是四维情形的粗估，严格的说法是紧接着的那句：作用量被 $I\cdot\mathrm{Vol}(\Sigma)$ 从下界住，其中 $I$ 是该解「每单位 $\mathbb{R}^{n-2}$ 体积」的作用量。$n=4$ 时 $\Sigma$ 是二维曲面、$\mathrm{Vol}(\Sigma)\sim R^2$，而涡旋（vortex）每单位面积的作用量具有质量平方量纲、量级为 $m^2$（相差 $1/e^2$ 型因子），于是 $I\cdot\mathrm{Vol}(\Sigma)\sim(mR)^2$；一般 $n$ 下则是 $I\,R^{n-2}$。这条下界只依赖 $c_1(\mathcal{L})$ 的 Poincaré 对偶圈、与 $M$ 的具体形状无关，这正是下文「面积律只由 $C$ 张成的极小曲面决定」的伏笔。]

的确，若 $\phi$ 是 $\mathcal{L}$ 的截面，则 $\phi$ 必须在与 $c_1(\mathcal{L})$ 对偶的余维二维圈 $\Sigma$ 上取零。
固定 $\Sigma$ 后，我们可以考察带这种零点的最小作用量组态。
为此，我们可以（若把 $M$ 的度量放大）化归为 $M=\mathbb{R}^n=\mathbb{R}^{n-2}\times\mathbb{R}^2$、$\Sigma$ 等于第一个因子的情形。
我们可以假定 $\phi$ 与 $A$ 在 $M$ 的第一个因子的平移下不变。
在第二个因子中，我们要求 $\phi$ 在原点取零，并在无穷远处（在规范变换意义下）趋于真空，使得丛的第一 Chern 类——相对于 $\phi$ 在无穷远处给出的平凡化——等于 1。
同样的问题在第 7.3 讲出现过（以寻找同一模型二维版本的瞬子解的形式），

% ===== 原文 p.142 =====
而我们讨论过该解的定性性质。
总之，设 $I$ 为该解在二维意义下的作用量（即只对 $M$ 中第二个因体积分）。
那么在 $n$ 维意义下，$I$ 是该解每单位体积 $\mathbb{R}^{n-2}$ 的作用量。
回到流形 $M$ 上的整体问题，设 $\Sigma\subset M$ 为代表 $\mathcal{L}$ 的第一 Chern 类的、体积最小的余维二维圈。
最小作用量场在 $\Sigma$ 的法向上看起来就像刚才描述的瞬子；其作用量近似为 $I\cdot\mathrm{Vol}(\Sigma)$。

现在让我们特殊化到四维，并考虑 $M=\mathbb{R}^4$ 中一个圆 $C$ 上的 't Hooft 圈。\footnote{我们其实可以在任意维数作这一讨论，但在 $n$ 维中一般地，'t Hooft 算符定义在一个余维三的闭子流形上；只有在 $n=4$ 时它才是圈算符。然而，涉及 't Hooft 圈与 Wilson 圈对偶的更深层陈述是 $n=4$ 特有的。}
在此情形，规范丛只定义在 $M\setminus C$ 上。
若 $\mathcal{L}$ 是 $M\setminus C$ 上的线丛、它在连接 $C$ 的二维球面上 Chern 类为 1，而 $\phi$ 是 $\mathcal{L}$ 的截面，则 $\phi$ 必须在某个以 $C$ 为边界的二维曲面上取零。
这个零点将出现在张成 $C$ 的最小面积曲面 $D$ 上，而在 $D$ 的法向上该解看起来与上一段完全一样。
于是，与上面相同的论证表明 't Hooft 圈的期望值为 $\langle T_1(C)\rangle\sim e^{-I\,\mathrm{Area}(D)}$，其中 $\mathrm{area}(D)$ 是张成 $C$ 的曲面的最小面积。
这就是我们想要证明的面积律。

这种行为是所谓的 Higgs 区（或相）的特征。

现在让我们更详细地讨论已确立的 't Hooft 圈算符行为与禁闭的关系。
若 $\pi_1(G)\neq0$，我们迄今至少刻画了三种可能的相：
\begin{enumerate}
\item Coulomb：没有质量间隙，红外有质量为零的规范玻色子，$W$ 与 $T$ 的行为像自由理论，表现出长度律。
\item Higgs：有质量间隙，$W$ 满足长度律，$T$ 满足面积律。
\item 禁闭：有质量间隙，$W$ 满足面积律，$T$ 满足长度律。
\end{enumerate}
（由于多种原因，这并不是对相的完整分类。）

正如已经提示的，'t Hooft 的想法是：应当存在一种非 Abel 的对偶，它交换 $W$ 与 $T$，交换 Higgs 区与禁闭区，并把 Coulomb 相映到自身。
于是，一个理论中 $T$ 的面积律意味着对偶理论中的禁闭。
这正是某些超对称理论中所发生的事情；在 Donaldson 理论所涉及的那个理论中我们也会发现类似的东西。

事实上，'t Hooft 证明了：若 $R$ 与 $\gamma$ 使得 $\gamma|_R\neq1$，则 $W_R$ 或 $T_\gamma$ 之中必有一个表现出面积律。
注意群 $\pi_1(G)^*\times\pi_1(G)$ 具有自然的辛形式（symplectic form）。
模去表示上的等价关系，Wilson 圈由第一个因子分类，'t Hooft 圈由第二个因子分类。
给定一对 $(c,\gamma)\in\pi_1(G)^*\times\pi_1(G)$，取具有中心特征标（central character）$c$ 的表示 $R$，并考虑算符 $W_R(C)T_\gamma(C)$（它可能需要一些重整化，但这不影响下面的评论）。
't Hooft 证明了：使得该算符不表现出面积律的那些 $(c,\gamma)$ 的集合，是 $\pi_1(G)^*\times\pi_1(G)$ 关于自然辛形式的一个迷向子群（isotropic subgroup）$H$。
其理由大致如下。
若 $A(C)=W_{R_1}T_{\gamma_1}(C)$ 与 $B(C)=W_{R_2}T_{\gamma_2}(C)$ 表现出长度律，则作用在真空上时，它们只产生沿 $C$ 局域化的效应（否则就会有面积律，

% [Tutor Note: 「迷向」可以逐字用 (10.4) 验证，不必凭印象：令 $A=W_{R_1}T_{\gamma_1}$、$B=W_{R_2}T_{\gamma_2}$，用 (10.4) 把 $T$ 逐个移到右边（$W$ 之间、$T$ 之间各自对易），得 $AB=q\,BA$，其中 $q=\exp\bigl(2\pi i(\langle c_2,\gamma_1\rangle-\langle c_1,\gamma_2\rangle)\bigr)=\exp(-2\pi i\,\omega)$；这里 $\langle c,\gamma\rangle$ 是 $\pi_1(G)^*$ 与 $\pi_1(G)$ 的配对（即 $R$ 的中心特征标在 $\gamma$ 上的取值），$\omega\bigl((c,\gamma),(c',\gamma')\bigr)=\langle c,\gamma'\rangle-\langle c',\gamma\rangle$。于是 $q=1$ 恰等价于 $H$ 关于 $\omega$ 迷向。对照本章对三个相的列举：Higgs 相对应 $H=\pi_1(G)^*\times\{0\}$（所有 $W$ 都走长度律），禁闭相对应 $H=\{0\}\times\pi_1(G)$（所有 $T$ 都走长度律），二者都是极大迷向子群——这正是作者说「每个有质量相关联着一个极大迷向子群」的含义。]

% ===== 原文 p.143 =====
而不是长度律）。所以
\begin{equation}\tag{10.11}\label{eq:10-11}
\langle A(C)B(C')\rangle=\langle B(C')A(C)\rangle=\langle A(C)\rangle\langle B(C')\rangle\bigl(1+o(1)\bigr),\qquad d\to\infty
\end{equation}
其中 $d$ 是 $C$ 与 $C'$ 之间的距离。
这表明：若 $AB=qBA$，其中 $q$ 是常数，则 $q$ 必须等于 1。
由 't Hooft 的公式 (10.4)（它在非 Abel 情形显然也成立），这意味着 $H$ 是迷向的。
进一步的物理论证表明：若有质量间隙，则 $H$ 是极大的迷向子群。
这导向比上面所述更精细的有质量相的分类：每一个有质量相都关联着 $\pi_1(G)^*\times\pi_1(G)$ 的一个极大迷向子群。
一般地，所有可能性都可以出现。

% [Tutor Note: (10.11) 是集团分解（cluster decomposition）在大间隔下的写法，右端的 $1+o(1)$ 正是质量间隙的定量内容：修正项应为 $O(e^{-md})$（$m$ 为最低质量），若没有间隙，修正只会是幂次小，$d\to\infty$ 时不足以断言 $q=1$。论证的链条是：长度律 $\Rightarrow$ 两个算符的效应各自局域在自己的圈附近（相距 $d$）$\Rightarrow$ 大 $d$ 时因子化 $\Rightarrow$ $q=1$ $\Rightarrow$ 由 (10.4) 得 $\omega|_H=0$。下面的 Remark 1 恰恰指出这条链条在面积律处断裂：此时效应局域在极小曲面 $D$ 上，而 $D$ 必与环绕数非零的 $C'$ 相交，$d$ 恒为 0，因子化无从谈起。]

\textbf{注 1（Remark 1）.}
证明圈算符 $A$ 与 $B$ 必须对易的论证，在其中一个算符（比如 $A$）表现出面积律时失效。
此时，该算符作用在真空上产生的效应完全不是局域在 $C$ 附近的；它产生的效应局域在以 $C$ 为边界的极小面积圆盘 $D$ 附近。
当 $C$ 与 $C'$ 的环绕数不为零时，这样的圆盘总会与 $C'$ 相交。
公式 (10.11) 现在只在 $d$ 是 $C'$ 到 $D$ 的距离时才成立，而该距离总是 0，所以公式不能告诉我们任何东西。

\textbf{注 2（Remark 2）.}
Higgs 相中 't Hooft 算符的面积律有许多物理与数学上的应用。
例如，稍作调整之后，我们上面关于「带非零第一 Chern 类的丛的 Higgs 相行为」的分析，可以用作 Meissner 效应（Meissner effect）的一种解释，即超导体（近似由我们所考察的 Abel Higgs 模型描述）排斥磁通量这一事实。
读者也许在科学博览会上见过磁体悬浮在超导体之上的演示；这一效应有相同的起源。
在数学上，C. Taubes 对辛四维流形的 Seiberg–Witten 不变量的分析用到了同样的事实：在（一个密切类比的）方程组中，$\phi$ 的零点局域在面积最小的曲面上。

\textbf{参考文献与交叉引用（References and cross references）}

Abel 对偶在第 8 讲讨论。
Wilson 圈算符在第 7 讲定义。
§10.3 的理论在第 2 讲讨论。
四维中的非 Abel 对偶将在后面关于四维超对称场论的各讲中讨论。

关于 't Hooft 的禁闭图像，\textit{inter alia} 参见
\begin{quote}
G. 't Hooft, \textit{Aspects of quark confinement}, Physica Scripta \textbf{24} (1981), 841--846.\\
G. 't Hooft, \textit{The topological mechanism for permanent quark confinement in a non-abelian gauge theory}, Physica Scripta \textbf{25} (1982), 133--142.
\end{quote}
